\section{Тема 9. Линейная алгебра и сжатие информации}

\subsection{Постановка задачи (вариант~4)}

В матрице $A(m,n)$ выбран произвольный ненулевой минор $k$-го
порядка ($k \le m$, $k \le n$): $\det\,(a_{i_pj_q})$, $1\le p,q\le
k$. Минор задаётся множеством номеров строк и множеством номеров
столбцов мощности~$k$. Вычислить значение этого минора.

\subsection{Математическое обоснование}

Минор $k$-го порядка матрицы $A$ --- определитель квадратной
подматрицы, образованной выбранными строками $\{i_1,\ldots,i_k\}$
и столбцами $\{j_1,\ldots,j_k\}$. Определитель вычисляется
рекурсивно разложением по первой строке:
\[
  \det M = \sum_{c=0}^{k-1} (-1)^{c} \, m_{0c} \, \det M_{0c},
\]
где $M_{0c}$~--- подматрица без первой строки и столбца~$c$.
База рекурсии: $\det$ матрицы $1{\times}1$ равен её единственному
элементу, $2{\times}2$~--- $m_{00}m_{11}-m_{01}m_{10}$.

\subsection{Блок-схема алгоритма}

\begin{center}
\begin{tikzpicture}[
  node distance=10mm, start chain=going below,
  nodes={draw=black, top color=white, bottom color=lightgray!50,
         minimum width=5.6cm, align=center, font=\small},
  arrow/.style={->, >=Stealth, thick},
]
  \node[draw, rounded corners, on chain] {Start};
  \node[shape=trapezium, trapezium left angle=70,
        trapezium right angle=110, on chain, draw, join=by arrow]
       {$A(m,n)$,\ $k$,\ строки,\ столбцы};
  \node[shape=rectangle, on chain, draw, join=by arrow]
       {$M(k,k)$ из $A$};
  \node[shape=diamond, aspect=2, on chain, draw, join=by arrow,
        minimum width=5.2cm]
       (c1) {$k = 1$\,?};
  \node[shape=trapezium, trapezium left angle=70,
        trapezium right angle=110, draw, right=2cm of c1]
       (r1) {$m_{00}$};
  \draw[arrow] (c1.east) -- node[above,font=\scriptsize]{$+$} (r1.west);
  \node[shape=diamond, aspect=2, on chain, draw,
        below=14mm of c1, minimum width=5.2cm]
       (c2) {$k = 2$\,?};
  \draw[arrow] (c1.south) -- (c2.north);
  \node[shape=trapezium, trapezium left angle=70,
        trapezium right angle=110, draw, right=2cm of c2]
       (r2) {$m_{00}m_{11}{-}m_{01}m_{10}$};
  \draw[arrow] (c2.east) -- node[above,font=\scriptsize]{$+$} (r2.west);
  \node[shape=rectangle, on chain, draw, below=14mm of c2]
       (rec) {рекурсия по строке 0};
  \draw[arrow] (c2.south) -- (rec.north);
  \node[shape=trapezium, trapezium left angle=70,
        trapezium right angle=110, on chain, draw, join=by arrow]
       (res) {$\det(M)$};
  \node[draw, rounded corners, on chain, join=by arrow] (end) {End};
  \draw[arrow] (r1.south) |- (res.east);
  \draw[arrow] (r2.south) |- (res.east);
\end{tikzpicture}
\end{center}

\subsection{Текст программы}

\lstinputlisting[style=Cstyle]{../src/t09.c}

\subsection{Тестовые данные}

\begin{center}
\begin{tabular}{l|l}
\toprule
Минор & Значение \\
\midrule
$\begin{pmatrix}1&2\\3&4\end{pmatrix}$ & $-2$ \\
$\begin{pmatrix}1&0&2\\-1&3&1\\2&1&0\end{pmatrix}$ & $-5$ \\
\bottomrule
\end{tabular}
\end{center}

\textbf{Результат выполнения программы} (минор $2{\times}2$):
\begin{lstlisting}[language={},basicstyle=\ttfamily\small,frame=single]
Размеры матрицы M N: 2 2
Введите матрицу A(2x2):
1 2
3 4
Порядок минора k (<=4): 2
Введите 2 номеров строк: 0 1
Введите 2 номеров столбцов: 0 1
Значение минора: -2.0000
\end{lstlisting}
